\documentclass[11pt]{article}

\usepackage[a4paper,margin=27mm]{geometry}
\usepackage{amsmath,amssymb,bm,mathtools}
\usepackage{booktabs}
\usepackage{microtype}
\usepackage[hidelinks]{hyperref}

\newcommand{\F}{\bm F}
\newcommand{\C}{\bm C}
\newcommand{\B}{\bm B}
\newcommand{\I}{\bm I}
\newcommand{\Pola}{\bm P}
\newcommand{\tauK}{\bm\tau}
\newcommand{\sig}{\bm\sigma}
\newcommand{\cof}{\operatorname{cof}}
\newcommand{\tr}{\operatorname{tr}}
\newcommand{\dev}{\operatorname{dev}}

\title{An Inversion-Tolerant Nearly Incompressible\\
Mooney--Rivlin Formulation}
\author{}
\date{}

\begin{document}
\maketitle

\begin{abstract}
We describe an inversion-tolerant continuation of the standard nearly
incompressible Mooney--Rivlin material. The admissible physical branch is left
unchanged for $J=\det\F\ge J_c>0$. Below $J_c$, the singular unimodular powers
$J^{-2/3}$ and $J^{-4/3}$ are replaced by smooth polynomial continuations.
The resulting strain-energy density is finite and twice differentiable through
collapsed and inverted states, so Newton-type optimization can evaluate and
cross configurations with $J\le0$ without imposing an inversion barrier. The
continuation is a numerical extension only; material calibration and the exact
unimodular Mooney--Rivlin response are preserved on the physical branch.
\end{abstract}

\section{Objective}

For finite-strain soft materials, a common nearly incompressible
Mooney--Rivlin model uses the unimodular invariants
\begin{equation}
  \bar I_1 = J^{-2/3} I_1,
  \qquad
  \bar I_2 = J^{-4/3} I_2,
  \qquad
  J=\det\F,
\end{equation}
where
\begin{equation}
  I_1=\tr\C,\qquad
  I_2=\frac{1}{2}\left((\tr\C)^2-\tr(\C^2)\right),
  \qquad
  \C=\F^T\F.
\end{equation}
The standard physical energy is
\begin{equation}
  W_{\mathrm{MR}}(\F)
  =
  C_{10}(\bar I_1-3)
  +
  C_{01}(\bar I_2-3)
  +
  U(J),
  \label{eq:physical_mr}
\end{equation}
with a volumetric penalty such as
\begin{equation}
  U(J)=\frac{\kappa}{2}(J-1)^2.
  \label{eq:volumetric}
\end{equation}
For large $\kappa$, this is nearly incompressible.

The difficulty is that $J^{-2/3}$ and $J^{-4/3}$ are singular at $J=0$ and
undefined as real-valued functions for $J<0$. A line search or continuation
method that starts from a tangled mesh therefore cannot evaluate the physical
energy along paths that pass through collapsed or inverted elements.

The goal is to keep \eqref{eq:physical_mr} exactly unchanged for
$J\ge J_c>0$, while replacing only the singular scalar powers below $J_c$ by a
finite $C^2$ extension.

\section{Modified Strain-Energy Function}

Let $m\in\{2/3,4/3\}$ and define $h_m(J)$ as a continuation of $J^{-m}$:
\begin{equation}
  h_m(J)=
  \begin{cases}
    J^{-m}, & J\ge J_c,\\
    p_m(J), & J_{\min}<J<J_c,\\
    p_m(J_{\min}), & J\le J_{\min},
  \end{cases}
  \qquad J_{\min}<0<J_c.
  \label{eq:hm_def}
\end{equation}
The modified Mooney--Rivlin strain-energy density is then
\begin{equation}
  \boxed{
  W_c(\F)
  =
  C_{10}\left(h_{2/3}(J) I_1 - 3\right)
  +
  C_{01}\left(h_{4/3}(J) I_2 - 3\right)
  +
  \frac{\kappa}{2}(J-1)^2.
  }
  \label{eq:continued_mr}
\end{equation}
For $J\ge J_c$, $h_{2/3}(J)=J^{-2/3}$ and
$h_{4/3}(J)=J^{-4/3}$, so \eqref{eq:continued_mr} is exactly the calibrated
unimodular Mooney--Rivlin model \eqref{eq:physical_mr}. For $J<J_c$, it is an
artificial continuation used only to make the nonlinear optimization problem
evaluable through inversion.

\section{\texorpdfstring{Quartic $C^2$ Continuation}{Quartic C2 Continuation}}

For each exponent $m$, write
\begin{equation}
  z=\frac{J-J_c}{\Delta},
  \qquad
  \Delta=J_c-J_{\min}>0,
\end{equation}
so that $z=0$ at $J_c$ and $z=-1$ at $J_{\min}$. Use
\begin{equation}
  p_m(J)=a_0+a_1z+a_2z^2+a_3z^3+a_4z^4.
\end{equation}
At $J_c$, match the physical power and its first two derivatives:
\begin{align}
  p_m(J_c) &= J_c^{-m},\\
  p_m'(J_c) &= -mJ_c^{-m-1},\\
  p_m''(J_c) &= m(m+1)J_c^{-m-2}.
\end{align}
At $J_{\min}$, impose flattening,
\begin{equation}
  p_m'(J_{\min})=0,
  \qquad
  p_m''(J_{\min})=0,
\end{equation}
so the constant branch for $J\le J_{\min}$ is also $C^2$.

The first coefficients are
\begin{align}
  a_0 &= J_c^{-m},\\
  a_1 &= -m\Delta J_c^{-m-1},\\
  a_2 &= \frac{1}{2}m(m+1)\Delta^2J_c^{-m-2}.
\end{align}
The two flattening conditions give
\begin{equation}
  \boxed{
  a_3=\frac{4}{3}a_2-a_1,
  \qquad
  a_4=\frac{1}{2}(a_2-a_1).
  }
  \label{eq:quartic_coefficients}
\end{equation}
Thus $h_m$, $h_m'$, and $h_m''$ are continuous at both $J_c$ and $J_{\min}$.

\section{Stress of the Continued Model}

The first Piola stress of \eqref{eq:continued_mr} is
\begin{equation}
  \Pola_c
  =
  C_{10}\left(h_1\frac{\partial I_1}{\partial\F}
  +h_1'I_1\frac{\partial J}{\partial\F}\right)
  +
  C_{01}\left(h_2\frac{\partial I_2}{\partial\F}
  +h_2'I_2\frac{\partial J}{\partial\F}\right)
  +
  \kappa(J-1)\frac{\partial J}{\partial\F},
  \label{eq:general_piola}
\end{equation}
where, for compactness,
\begin{equation}
  h_1=h_{2/3}(J),\qquad h_2=h_{4/3}(J).
\end{equation}
In three dimensions,
\begin{equation}
  \frac{\partial I_1}{\partial\F}=2\F,
  \qquad
  \frac{\partial I_2}{\partial\F}
  =
  2(I_1\F-\F\C),
  \qquad
  \frac{\partial J}{\partial\F}=\cof\F.
\end{equation}
Substitution gives
\begin{equation}
  \boxed{
  \Pola_c
  =
  2C_{10}h_1\F
  +
  2C_{01}h_2(I_1\F-\F\C)
  +
  \left(C_{10}h_1'I_1+C_{01}h_2'I_2+\kappa(J-1)\right)\cof\F.
  }
  \label{eq:continued_piola_3d}
\end{equation}
On $J\ge J_c$, this reduces exactly to the usual unimodular
Mooney--Rivlin stress because the continued powers and their derivatives match
the physical powers.

\section{Plane-Strain Prototype}

The accompanying prototype uses a two-dimensional triangular mesh with the
plane-strain embedding
\begin{equation}
  \F_3=
  \begin{bmatrix}
    F_{11} & F_{12} & 0\\
    F_{21} & F_{22} & 0\\
    0 & 0 & 1
  \end{bmatrix}.
\end{equation}
For this embedding,
\begin{equation}
  J=\det\F,\qquad
  I_1=\F:\F+1,\qquad
  I_2=\F:\F+J^2,
  \label{eq:plane_strain_invariants}
\end{equation}
where $\F$ is the in-plane $2\times2$ deformation gradient. With
$G=\cof\F$, the in-plane stress can be written
\begin{equation}
  \Pola_c
  =
  2(C_{10}h_1+C_{01}h_2)\F
  +
  \left[
    C_{10}h_1'I_1
    +
    C_{01}(2Jh_2+h_2'I_2)
    +
    \kappa(J-1)
  \right]G.
  \label{eq:plane_strain_piola}
\end{equation}
The consistent material tangent follows by differentiating
\eqref{eq:plane_strain_piola}; it requires $h_m$, $h_m'$, and $h_m''$ at each
quadrature point.

\section{Physical-Branch Consistency}

For $J\ge J_c$, the modified energy is identical to the physical energy:
\begin{equation}
  W_c(\F)=W_{\mathrm{MR}}(\F).
\end{equation}
Its gradient and Hessian also agree because the scalar continuations match the
first and second derivatives of the singular powers at $J_c$ and equal the
physical powers for all larger $J$. Consequently, any calibrated parameters
$C_{10}$, $C_{01}$, and $\kappa$ remain valid for converged states whose
determinants stay above $J_c$.

The isochoric part of the physical branch is scale invariant:
\begin{equation}
  W_{\mathrm{iso}}(a\F)=W_{\mathrm{iso}}(\F),\qquad a>0.
\end{equation}
Differentiating with respect to $a$ at $a=1$ gives
\begin{equation}
  \Pola_{\mathrm{iso}}:\F=0,
\end{equation}
and therefore the isochoric Kirchhoff stress is deviatoric:
\begin{equation}
  \tr(\tauK_{\mathrm{iso}})=0.
\end{equation}
The volumetric branch contributes
\begin{equation}
  \tauK_{\mathrm{vol}}=J U'(J)\I,
\end{equation}
so the standard deviatoric--volumetric split is preserved exactly on the
admissible physical branch.

\section{Optimization Use}

The continuation is intended for nonlinear solution from tangled initial
states. The prototype uses:
\begin{enumerate}
  \item continuation in $J_c$ from a less singular model to the target cutoff;
  \item an absolute-eigenvalue projected Newton Hessian;
  \item Armijo backtracking without an inversion barrier;
  \item no positional tether or reference-state penalty.
\end{enumerate}
The continuation removes the evaluability obstruction at $J=0$, but it does
not make the finite-element energy globally convex. Inverted configurations
can still converge to non-reference local minima, so the choice of $J_c$,
$J_{\min}$, and the continuation shape should be tested statistically over
many distorted initial states.

\section{Implementation Notes}

For production kernels, the continuation is local to each quadrature point:
\begin{enumerate}
  \item evaluate $J$ once;
  \item evaluate both continued powers and their derivatives together;
  \item use the physical branch for $J\ge J_c$ and polynomial arithmetic below
        $J_c$;
  \item avoid sign-based branches on $J$ other than the cutoff comparison;
  \item keep $h_m$, $h_m'$, and $h_m''$ available for the tangent.
\end{enumerate}
In vectorized CPU kernels, the continuation branch should be implemented with
straight-line polynomial evaluation on masked lanes or by grouping quadrature
points by branch when the cost of masking is significant.

\section{Conclusion}

The modified energy \eqref{eq:continued_mr} applies the same
inversion-tolerant idea to nearly incompressible Mooney--Rivlin elasticity:
the physical material is untouched for $J\ge J_c$, while a finite $C^2$
extension is supplied for collapsed and inverted states. This separates
material modeling from nonlinear-solver globalization. The continuation is
not a new physical law for $J<0$; it is a numerical mechanism that allows
Newton iterations and line searches to move through otherwise unevaluable
states.

\begin{thebibliography}{9}

\bibitem{Ogden1997}
R.~W. Ogden,
\newblock \emph{Non-Linear Elastic Deformations},
\newblock Dover, 1997.

\bibitem{Smith2018}
B.~Smith, F.~de Goes, and T.~Kim,
\newblock ``Stable Neo-Hookean Flesh Simulation,''
\newblock \emph{ACM Transactions on Graphics}, 37(2), Article 12, 2018.
\newblock DOI: \href{https://doi.org/10.1145/3180491}{10.1145/3180491}.

\bibitem{Chen2024}
H.~Chen, H.-T.~D. Liu, D.~I.~W. Levin, C.~Zheng, and A.~Jacobson,
\newblock ``Stabler Neo-Hookean Simulation: Absolute Eigenvalue Filtering for
Projected Newton,''
\newblock \emph{SIGGRAPH Conference Papers 2024}, 2024.
\newblock DOI: \href{https://doi.org/10.1145/3641519.3657433}
{10.1145/3641519.3657433}.

\end{thebibliography}

\end{document}
